\documentclass[10pt,a4paper]{amsart}
\usepackage[margin=19mm]{geometry}
\usepackage{amsmath,amssymb,mathtools}
\usepackage[scr=rsfs,cal=euler]{mathalfa}
\usepackage{xfrac}
\usepackage[unicode,hidelinks,bookmarks=false]{hyperref}
\newcommand{\ie}{i.e.\ }
\newcommand{\cf}{cf.\ }
\newcommand{\resp}{respectively\ }
\newcommand{\Subsection}{\subsection}
\providecommand{\nobreakdash}{\nobreak\hbox{-}}
\setlength{\emergencystretch}{2em}
\multlinegap0pt
\usepackage{float}
\usepackage{amssymb}
\usepackage{bm}
\usepackage{xcolor}
\newtheorem{thm}{Theorem}
\newcommand{\dv}{\upd_{\upv}\!}
\newcommand{\es}{\operatorname{S}}
\newcommand{\intp}{\,\,\intpsym}
\newcommand{\intpsym}{\,\, \setlength{\unitlength}{0.4mm}
	\begin{picture}(0, 0)(5, 5)%
	\put(0, 4){\line(1,0){5}} \put(5,4){\line(0,1){8}}
	\end{picture}\,\,}
\newcommand{\mbJ}{\mathbf{J}}
\newcommand{\mbn}{\mathbf{n}}
\newcommand{\mbu}{\mathbf{u}}
\newcommand{\mbv}{\mathbf{v}}
\newcommand{\mcL}{\mathcal{L}}
\newcommand{\upd}{\operatorname{d}}
\newcommand{\upv}{\operatorname{v}}
\title{The difference variational bicomplex and multisymplectic systems}
\author{Linyu Peng, Peter E. Hydon}
\date{Source: arXiv:2307.13935v2 (CC BY 4.0)}
\begin{document}
\maketitle

\noindent\emph{Source and licence.} The difference variational bicomplex and multisymplectic systems. Version arXiv:2307.13935v2 (CC BY 4.0), distributed by arXiv under the Creative Commons Attribution 4.0 International licence, http://creativecommons.org/licenses/by/4.0/, as recorded in the arXiv licence metadata for this exact version and shown on the abstract page https://arxiv.org/abs/2307.13935v2. Full licence notice: \texttt{licenses/arxiv-2307.13935-CC-BY-4.0.txt}.\par\smallskip

\noindent\textit{Editorial scope.} Counted unit: the article's thm thm (source0.tex lines 1597--1606). The article's complete original proof of it is delivered with it (source0.tex lines 1608--1661, one \texttt{proof} environment of 54 lines). No mathematics is written, repaired or re-derived: the statement and the proof are the article's own lines, in the article's own notation, and no retained line is substituted or re-flowed.  No further source-local context is needed: every cross-reference of every retained line resolves inside the two delivered spans.  Nothing was omitted from any retained span, so the package carries no omission record.  No retained line contains a citation, so the wrapper carries no bibliography and no bibliography entry.  No retained line contains a figure environment, an image-inclusion command or drawing code, so no image is delivered and the package declares no asset dependency.  Editorial formatting only, in addition: an AMS \texttt{amsart} preamble that restates the article's own declarations and macros used by the retained lines, namely the environments \texttt{thm} on separate counters, together with the standard packages those lines need.  The retained environments print in this document as Theorem 1 in order of appearance, whereas the article prints the same items with the numbering of its own full document.  The complete native source of the article as distributed in the arXiv e-print for this exact version is bundled unchanged as \texttt{sources/144800/source0.tex}.
\begin{thm}
For each $k=0,1,2,\ldots,p$, the vertical complex
\begin{equation*}
\Omega^{k,0}\stackrel{\dv }{\longrightarrow}
\Omega^{k,1}\stackrel{\dv }{\longrightarrow}
\Omega^{k,2}\stackrel{\dv }{\longrightarrow}
\cdots
\end{equation*}
is exact.
\end{thm}

\begin{proof}
Denote the prolongation of the vertical vector field
$u^{\alpha}\frac{\partial}{\partial u^{\alpha}}$ as
\begin{equation}
\mbv={(\es_{\mbJ}u^{\alpha})\frac{\partial}{\partial
u^{\alpha}_{\mbJ}}} ={u^{\alpha}_{\mbJ}\frac{\partial}{\partial
u^{\alpha}_{\mbJ}}},
\end{equation}
the flow generated by which on $\mathbb{Z}^p\times P(\mathbb{R}^q)$ is a
one-parameter family of diffeomorphisms
\begin{equation}
\exp(\varepsilon \mbv)(\mbn,\mbu)=(\mbn,\ldots,e^{\varepsilon}u^{\alpha}_{\mbJ},\ldots),
\end{equation}
satisfying\footnote{ The mapping $\exp$ is  called the
exponential mapping, and we claim that it is a diffeomorphism here
as the discrete parts can be considered as fixed parameters.}
\begin{equation}
\frac{\upd}{\upd\!\varepsilon}\Big|_{\varepsilon=0}\exp(\varepsilon
\mbv)(\mbn,\mbu)=\mbv.
\end{equation}
 As a consequence, naturally there exist induced push-forward and
pull-back mappings. Take any form $\sigma\in\Omega^{k,l}$. As the
transformation only occurs in the continuous parts, we can define
a derivative of $\sigma$ as\footnote{ In the continuous setting, this derivative
is exactly the Lie derivative, and we also refer to its notation and
call it the Lie derivative accordingly. It satisfies all the
properties that the canonical Lie derivative owns.}
\begin{equation}\label{lieddiscrete}
\mcL_{\mbv}\sigma=\lim_{\varepsilon\rightarrow0}\frac{\exp(\varepsilon \mbv)^{\ast}(\sigma)-\sigma}{\varepsilon}=\frac{\upd}{\upd\!\varepsilon}\Big|_{\varepsilon=0}
{\exp(\varepsilon \mbv)^{\ast}(\sigma)},
\end{equation}
which implies that
\begin{equation}
\frac{\upd}{\upd\!\varepsilon}{\exp(\varepsilon \mbv)^{\ast}(\sigma)}=
\exp(\varepsilon \mbv)^{\ast}(\mcL_\mbv\sigma).
\end{equation}

The
vertical homotopy operators
$h_{\upv}^{k,l}:\Omega^{k,l}\rightarrow\Omega^{k,l-1},l\geq1$ are defined by
\begin{equation}\label{eq:vehomo}
h^{k,l}_{\upv}(\sigma)=\int_0^1{\frac{1}{\lambda}\exp(\ln\lambda
\mbv)^{\ast}(\mbv\intp\sigma) \upd\!\lambda},
\end{equation}
such that
\begin{equation}
\begin{aligned}
\sigma=\dv\,(h_{\upv}^{k,l}(\sigma))
+h_{\upv}^{k,l+1}(\dv \sigma).
\end{aligned}
\end{equation}
The integrand in \eqref{eq:vehomo} is a smooth function of
$\lambda$ at $\lambda=0$, as $l\geq1$. This finishes the proof.
\end{proof}

\end{document}
